\chapter{Topological Fixed Point Theorems}

\section*{11.1. The Schauder Fixed Point Theorem}

\begin{theorem}[Schauder fixed point theorem]
\label{gt-thm-11-1}
\dcref{ch11:11.1}
\group{gt-ch11}
\source{gilbarg-trudinger:page-0291,gilbarg-trudinger:page-0292}
Let $\G$ be a compact convex subset of a Banach space $\B$. Every continuous
map $T:\G\to\G$ has a fixed point.
\end{theorem}
\begin{proof}
\sketch Cover $\G$ by finitely many balls of radius $1/k$, take the convex
hull $\G_k$ of their centers, and use distance-weighted barycentric
coordinates to construct a continuous map $J_k:\G\to\G_k$ satisfying
\begin{equation}
\tag{11.1}
\|J_kx-x\|<\frac1k.
\end{equation}
Brouwer's theorem gives a fixed point $x_k$ of $J_kT$ on $\G_k$.
Compactness supplies a convergent subsequence, and (11.1) implies that its
limit is fixed by $T$.
\end{proof}

\begin{corollary}[Schauder fixed point theorem for compact maps]
\label{gt-cor-11-2}
\source{gilbarg-trudinger:page-0292}
\dcref{ch11:11.2}
\group{gt-ch11}

Let $\G$ be closed and convex in a Banach space. If $T:\G\to\G$ is
continuous and $T(\G)$ is precompact, then $T$ has a fixed point.
\end{corollary}
\begin{proof}
\sketch Apply Theorem~\ref{gt-thm-11-1} to the closed convex hull of
$T(\G)$, which is a compact subset of $\G$ invariant under $T$.
\end{proof}

\section*{11.2. A Leray--Schauder Principle}

\begin{theorem}[Leray--Schauder principle]
\label{gt-thm-11-3}
\source{gilbarg-trudinger:page-0292}
\uses{gt-cor-11-2}
\dcref{ch11:11.3}
\group{gt-ch11}


Let $T:\B\to\B$ be compact. Suppose that some $M$ satisfies
\begin{equation}
\tag{11.2}
\|x\|_{\B}<M
\end{equation}
whenever $x=\sigma Tx$ for $\sigma\in[0,1]$. Then $T$ has a fixed point.
\end{theorem}
\begin{proof}
\uses{gt-cor-11-2}
\sketch After scaling, take $M=1$ and radially retract $T$ to a compact map
of the closed unit ball into itself. Corollary~\ref{gt-cor-11-2} supplies a
fixed point $x$ of the retracted map. If $\|Tx\|\geq1$, then
$x=\sigma Tx$ with $\|x\|=1$ for some $\sigma\in(0,1]$, contrary to (11.2).
Hence $x=Tx$.
\end{proof}

Fix $0<\beta<1$ and consider
\begin{equation}
\tag{11.3}
Qu=a^{ij}(x,u,Du)D_{ij}u+b(x,u,Du)
\end{equation}
on a bounded domain $\Omega$. For $v\in C^{1,\beta}(\overline\Omega)$,
let $Tv$ solve
\begin{equation}
\tag{11.4}
a^{ij}(x,v,Dv)D_{ij}(Tv)+b(x,v,Dv)=0,
\qquad Tv|_{\partial\Omega}=\varphi.
\end{equation}
Then $u=Tu$ is equivalent to the Dirichlet problem for $Q$, while
$u=\sigma Tu$ is equivalent to
\begin{equation}
\tag{11.5}
Q_\sigma u=a^{ij}(x,u,Du)D_{ij}u+\sigma b(x,u,Du)=0,
\qquad u|_{\partial\Omega}=\sigma\varphi.
\end{equation}

\begin{theorem}[A priori criterion for Dirichlet solvability]
\label{gt-thm-11-4}
\source{gilbarg-trudinger:page-0293,gilbarg-trudinger:page-0294}
\uses{gt-thm-11-3}
\dcref{ch11:11.4}
\group{gt-ch11}


Let $\partial\Omega\in C^{2,\alpha}$,
$\varphi\in C^{2,\alpha}(\overline\Omega)$, and suppose $Q$ is elliptic on
$\overline\Omega$ with $a^{ij},b\in C^1$ on
$\overline\Omega\times\mathbb R\times\mathbb R^n$. If for some $\beta>0$
every $C^{2,\alpha}$ solution of (11.5), uniformly in $\sigma\in[0,1]$,
satisfies
\begin{equation}
\tag{11.6}
\|u\|_{C^{1,\beta}(\overline\Omega)}<M,
\end{equation}
then $Qu=0$ with $u=\varphi$ on $\partial\Omega$ has a solution in
$C^{2,\alpha}(\overline\Omega)$.
\end{theorem}
\begin{proof}
\uses{gt-thm-11-4}
\sketch Linear Schauder theory makes the operator in (11.4) well defined.
The global estimate maps bounded subsets of $C^{1,\beta}$ to bounded subsets
of $C^{2,\alpha\beta}$, hence to precompact subsets of $C^{1,\beta}$.
Passing to the limit in (11.4) proves continuity. Theorem~\ref{gt-thm-11-3}
and (11.6) therefore give a fixed point; Schauder regularity improves it to
$C^{2,\alpha}$.
\end{proof}

\section*{11.3. A Bounded-Slope Application}

The boundary graph $(\partial\Omega,\varphi)$ satisfies the
\emph{bounded-slope condition} with constant $K$ if at each
$P\in(\partial\Omega,\varphi)$ there are affine functions $\pi_P^\pm$ through
$P$ such that
\[
\pi_P^-\leq\varphi\leq\pi_P^+\quad\hbox{on }\partial\Omega,
\qquad |D\pi_P^\pm|\leq K.
\]
Consider either
\begin{equation}
\tag{11.7}
Qu=\operatorname{div}A(Du)
\end{equation}
or, in dimension two,
\begin{equation}
\tag{11.8}
Qu=a^{ij}(x,u,Du)D_{ij}u.
\end{equation}

\begin{theorem}[Existence under the bounded-slope condition]
\label{gt-thm-11-5}
\dcref{ch11:11.5}
\group{gt-ch11}
\source{gilbarg-trudinger:page-0295,gilbarg-trudinger:page-0296,gilbarg-trudinger:page-0297}
Assume that $Q$, $\Omega$, and $\varphi$ satisfy the hypotheses of
Theorem~\ref{gt-thm-11-4}, that $Q$ has form (11.7) or (11.8), and that
$(\partial\Omega,\varphi)$ satisfies the bounded-slope condition. Then the
Dirichlet problem $Qu=0$, $u=\varphi$ on $\partial\Omega$, has a solution in
$C^2(\overline\Omega)$.
\end{theorem}
\begin{proof}
\sketch For solutions with boundary value $\sigma\varphi$, the maximum
principle and the affine barriers give
\begin{equation}
\tag{11.9}
\sup_\Omega|u|=\sigma\sup_{\partial\Omega}|\varphi|
\leq\sup_{\partial\Omega}|\varphi|,
\end{equation}
and
\begin{equation}
\tag{11.10}
\sup_{\partial\Omega}|Du|\leq\sigma K\leq K.
\end{equation}
For (11.7), the weak equation is
\begin{equation}
\tag{11.11}
\int_\Omega A(Du)\cdot D\eta\,dx=0
\qquad(\eta\in C_0^1(\Omega)).
\end{equation}
Differentiating weakly shows that $w=D_ku$ satisfies
\begin{equation}
\tag{11.12}
D_i(\bar a^{ij}(x)D_jw)=0,
\qquad \bar a^{ij}(x)=a^{ij}(Du(x)).
\end{equation}
In dimension two, dividing (11.8) by one principal coefficient gives
analogously
\begin{equation}
\tag{11.14}
D_i(a_1^{ij}(x)D_jw)=0,
\qquad i,j=1,2.
\end{equation}
The weak maximum principle therefore yields
\begin{equation}
\tag{11.13}
\sup_\Omega|Du|=\sup_{\partial\Omega}|Du|\leq K.
\end{equation}
Interior Harnack estimates for these derivative equations, boundary
flattening, the boundary Harnack estimate for tangential derivatives, and the
equation for the normal derivative first give, for
$\Omega'\Subset\Omega$ and $d=\dist(\Omega',\partial\Omega)$,
\begin{equation}
\tag{11.15}
[Du]_{\beta;\Omega'}\leq Cd^{-\beta},
\end{equation}
and then give
\begin{equation}
\tag{11.16}
[Du]_{\beta;\Omega}\leq C
\end{equation}
uniformly in $\sigma$. Theorem~\ref{gt-thm-11-4} applies.
\end{proof}

\section*{11.4. The Leray--Schauder Fixed Point Theorem}

\begin{theorem}[Parameterized Leray--Schauder theorem]
\label{gt-thm-11-6}
\dcref{ch11:11.6}
\group{gt-ch11}
\source{gilbarg-trudinger:page-0298,gilbarg-trudinger:page-0299}
Let $T:\B\times[0,1]\to\B$ be compact and satisfy $T(x,0)=0$. If some
$M$ satisfies
\begin{equation}
\tag{11.17}
\|x\|_{\B}<M
\end{equation}
whenever $x=T(x,\sigma)$, then $T_1(x)=T(x,1)$ has a fixed point.
\end{theorem}
\begin{proof}
\sketch Normalize $M=1$. On the closed unit ball, interpolate between the
maps $x\mapsto T(x/(1-\varepsilon),1)$ and
$x\mapsto T(x/\|x\|,(1-\|x\|)/\varepsilon)$ near the boundary. The resulting
compact map sends the boundary to $0$, so the radial retraction argument has
a fixed point $x_\varepsilon$. For $\varepsilon=1/k$, compactness gives a
limit $(x,\sigma)$. If $\sigma<1$, then $\|x\|=1$ and
$x=T(x,\sigma)$, contradicting (11.17). Thus $\sigma=1$ and $x=T(x,1)$.
\end{proof}

\begin{lemma}[Inward boundary fixed-point criterion]
\label{gt-lem-11-7}
\dcref{ch11:11.7}
\group{gt-ch11}
\source{gilbarg-trudinger:page-0298}
Let $B$ be the open unit ball in $\B$. If
$T:\overline B\to\B$ is continuous, $T(\overline B)$ is precompact, and
$T(\partial B)\subset B$, then $T$ has a fixed point.
\end{lemma}
\begin{proof}
\sketch Radially retract $T$ to $\overline B$ and apply
Corollary~\ref{gt-cor-11-2}. A fixed point on the sphere would contradict
$T(\partial B)\subset B$, so the retraction is inactive at the fixed point.
\end{proof}

For a family
\[
Q_\sigma u=a^{ij}(x,u,Du;\sigma)D_{ij}u+b(x,u,Du;\sigma),
\]
assume:
\begin{enumerate}
\item $Q_1=Q$ and $b(x,z,p;0)=0$;
\item each $Q_\sigma$ is elliptic;
\item $a^{ij}$ and $b$ depend continuously on $\sigma$ as
$C^\alpha(\overline\Omega\times\mathbb R\times\mathbb R^n)$ functions.
\end{enumerate}

\begin{theorem}[Continuation theorem]
\label{gt-thm-11-8}
\dcref{ch11:11.8}
\group{gt-ch11}
\source{gilbarg-trudinger:page-0299}
Let $\partial\Omega\in C^{2,\alpha}$ and
$\varphi\in C^{2,\alpha}(\overline\Omega)$. Suppose the family $Q_\sigma$
satisfies the preceding conditions and, for some $\beta>0$, every solution of
\[
Q_\sigma u=0\quad\hbox{in }\Omega,
\qquad u=\sigma\varphi\quad\hbox{on }\partial\Omega
\]
satisfies $\|u\|_{C^{1,\beta}(\overline\Omega)}<M$, uniformly in
$\sigma\in[0,1]$. Then $Qu=0$, $u=\varphi$ on $\partial\Omega$, has a
solution in $C^{2,\alpha}(\overline\Omega)$.
\end{theorem}
\begin{proof}
\sketch Freeze the coefficients at $v$ and solve the associated linear
Dirichlet problem to obtain a compact continuous map
$T(v,\sigma)$ on $C^{1,\beta}(\overline\Omega)$. The assumptions give
$T(v,0)=0$, identify fixed points with the displayed family, and provide the
uniform bound required by Theorem~\ref{gt-thm-11-6}.
\end{proof}

\section*{11.5. Variational Problems}

Let $\Omega$ be bounded, let $F\in C^1(\Omega\times\mathbb R\times\mathbb
R^n)$, and set
\begin{equation}
\tag{11.18}
I(u)=\int_\Omega F(x,u,Du)\,dx.
\end{equation}
For fixed $\varphi\in C^{0,1}(\overline\Omega)$, let
\[
\mathcal C=\{u\in C^{0,1}(\overline\Omega):u=\varphi
\text{ on }\partial\Omega\}.
\]
Problem $\mathcal P$ asks for a minimizer of $I$ on $\mathcal C$. A minimizer
satisfies
\begin{equation}
\tag{11.19}
\int_\Omega\bigl(D_{p_i}F(x,u,Du)D_i\eta+D_zF(x,u,Du)\eta\bigr)\,dx=0,
\end{equation}
equivalently the weak Euler--Lagrange equation
\begin{equation}
\tag{11.20}
Qu=\operatorname{div}D_pF(x,u,Du)-D_zF(x,u,Du)=0.
\end{equation}
Call $I$ \emph{regular} when $F$ is strictly convex in $p$.

\begin{theorem}[Variational and Euler--Lagrange solvability]
\label{gt-thm-11-9}
\dcref{ch11:11.9}
\group{gt-ch11}
\source{gilbarg-trudinger:page-0301}
If $I$ is regular and $F$ is jointly convex in $(z,p)$, problem $\mathcal P$
has at most one solution. Its solvability is equivalent to solvability in
$C^{0,1}(\overline\Omega)$ of (11.20) with boundary value $\varphi$.
\end{theorem}
\begin{proof}
\sketch A minimizer has zero first variation, hence solves (11.19).
Conversely, if $u$ solves (11.19), convexity makes
$t\mapsto I(u+t(v-u))$ convex with derivative zero at $t=0$, so
$I(u)\leq I(v)$. Strict convexity in the gradient and the comparison
principle give uniqueness.
\end{proof}

For $K>0$, put
\[
\mathcal C_K=\{u\in\mathcal C:\|u\|_{C^{0,1}(\Omega)}\leq K\}.
\]
Problem $\mathcal P_K$ minimizes $I$ over $\mathcal C_K$; its minimizers are
called $K$-quasisolutions of $\mathcal P$.

\begin{theorem}[Existence of quasisolutions]
\label{gt-thm-11-10}
\dcref{ch11:11.10}
\group{gt-ch11}
\source{gilbarg-trudinger:page-0302,gilbarg-trudinger:page-0303}
Suppose $F\in C^1(\Omega\times\mathbb R\times\mathbb R^n)$, the functions
$F,D_zF,D_{p_i}F$ extend continuously to
$\overline\Omega\times\mathbb R\times\mathbb R^n$, and $F$ is convex in
$p$. If $\mathcal C_K$ is nonempty, then $\mathcal P_K$ has a solution.
\end{theorem}
\begin{proof}
\sketch The set $\mathcal C_K$ is precompact for uniform convergence.
Convexity in $p$ gives, for $u_m\to u$ uniformly,
\[
I(u_m)-I(u)\geq-o(1)+
\int_\Omega D_pF(x,u,Du)\mathbin{\cdot}D(u_m-u)\,dx.
\]
Approximate the bounded coefficient $D_pF(x,u,Du)$ in $L^1$ by compactly
supported $C^1$ fields and integrate by parts. The last term tends to zero,
so $I$ is lower semicontinuous. A minimizing sequence therefore has a
uniformly convergent subsequence whose limit minimizes $I$.
\end{proof}

\begin{theorem}[An interior quasisolution is a minimizer]
\label{gt-thm-11-11}
\dcref{ch11:11.11}
\group{gt-ch11}
\source{gilbarg-trudinger:page-0303,gilbarg-trudinger:page-0304}
Let $u$ be a $K$-quasisolution satisfying
\begin{equation}
\tag{11.22}
\|u\|_{C^{0,1}(\overline\Omega)}<K.
\end{equation}
If $F$ is jointly convex in $(z,p)$, then $u$ solves $\mathcal P$.
\end{theorem}
\begin{proof}
\sketch For $v\in\mathcal C$, condition (11.22) permits
$w=(1-\varepsilon)u+\varepsilon v\in\mathcal C_K$ for small
$\varepsilon>0$. Minimality on $\mathcal C_K$ and convexity give
\[
I(u)\leq I(w)\leq(1-\varepsilon)I(u)+\varepsilon I(v),
\]
hence $I(u)\leq I(v)$.
\end{proof}
